%% notes-tex/019-simb-multimodal/sections/9-campaign.tex
%% SOURCE: cost figures from experiments/019-simb-multimodal/results/
%% expression_objective_diagnosis.json (measured _runtime and epoch counts) and
%% round_leaderboards.csv. The design is proposed, not measured, and says so.

\section{The expression and morphology campaign\secstatus{todo}}
\label{sec:campaign}

The next campaign is expression plus morphology, for Figure 3. What follows is sized in
GPU-days rather than in arms, because the measured cost per run is what makes most possible
designs unaffordable.

\subsection{The budget arithmetic, measured\secstatus{todo}}

\begin{table}[htbp]
\centering\small
%% SOURCE: results/expression_objective_diagnosis.json (_runtime, epoch)
\begin{tabular}{lrrrr}
\toprule
run & epochs & s/epoch & wall clock & peak at \\
\midrule
v9 \file{M_fine} (masked)  & 9{,}999 & 32.9 & \textbf{91.4 h = 3.81 d} & epoch 9{,}674 ($97\,\%$) \\
v8 \file{V_basis64}        & 4{,}106 & 42.0 & 47.9 h = 2.00 d & epoch 3{,}921 ($95\,\%$) \\
\bottomrule
\end{tabular}
\caption[]{One expression run is a multi-day object. Both runs peaked in their last few
percent, so neither bracketed its own maximum.}
\label{tab:campaign-cost}
\end{table}

At $\approx 35$ s/epoch, one GPU delivers $\approx 2{,}500$ epochs per day. So with $G$
GPUs for $D$ days the campaign has

$$\text{arm-slots} \;=\; \frac{G \cdot D \cdot 2{,}500}{E}$$

runs, where $E$ is the per-arm epoch budget. \textbf{Every design decision below is a
consequence of that fraction}, so it should be evaluated with the real $G$ and $D$ before
anything is launched.

\subsection{Set the epoch budget at 4,000, and the reason is measurable\secstatus{todo}}

The previous round's failure was scoring arms at 300 epochs against a curve that had not
turned. The opposite failure is affordable only once: 10,000 epochs is 3.8 days per arm,
which at four GPUs for four days buys four runs and no replicates.

$E = 4{,}000$ is the defensible middle, and not by taste. The v8 run \textbf{peaked at
epoch 3,921} of 4,106, so a pair-rank arm has been observed to reach its maximum inside
that budget. It costs $\approx 1.6$ days per run, so four GPUs for four days buys
$\approx 10$ slots.

\notesec{What $E = 4{,}000$ gives up, stated plainly} The v9 masked run was still climbing
at 9,674 and reached $0.2362$; truncating it at 4,000 would have scored it near $0.21$.
So a 4,000-epoch comparison is fair between arms and \textbf{understates} the absolute
number by roughly $0.02$ to $0.03$. That is the right trade for an arm comparison and the
wrong one for a headline figure, which is why the winning arm gets one long run afterward.

\subsection{Seeds: two per arm, and the noise floor is already known\secstatus{todo}}

Expression's across-initialization spread was measured at $\approx 0.006$ ($n = 8$, fixed
split) in wave 5, an order of magnitude below the betaxanthin strand's $0.030$
(\cref{tab:replicate-noise}), because the split is pinned so \file{seed} varies weights
only. At $\sigma = 0.006$, two seeds give a paired standard error of $\approx 0.004$, so a
$0.015$ effect is nearly $4\sigma$. \textbf{Two seeds per arm is enough}; three would buy
precision the mechanism differences do not need.

\subsection{Phase A, the objective, and why it goes first\secstatus{todo}}

\cref{sec:objective} found that the quantile loss reaches its minimum at epoch 463 and
rises for the next 9,500 while Pearson climbs, and that the model never gets \file{nmse}
below 1. A mechanism arm evaluated under an objective that is optimizing something else
cannot separate ``the mechanism does not help'' from ``the objective does not reward it'',
so the objective is settled first.

\begin{table}[htbp]
\centering\small
\begin{tabular}{llp{62mm}}
\toprule
arm & objective & what a win means \\
\midrule
\texttt{O\_ref}   & quantile (current) & reference \\
\texttt{O\_zmse}  & MSE on per-gene $z$-scored targets & ``predict the mean'' scores 0 rather than winning; already implemented as \file{standardize_per_feature_target} \\
\texttt{O\_corr}  & correlation objective & optimizing the quantity actually scored \\
\bottomrule
\end{tabular}
\caption[]{Phase A. Three arms, two seeds, $E = 4{,}000$: six slots, $\approx 10$ GPU-days.}
\label{tab:phase-a}
\end{table}

\notesec{Report \file{nmse} and the SD ratio on every arm, not just Pearson} The
calibration identity in \cref{sec:objective} means an arm can raise Pearson while leaving
\file{nmse} above 1, and that is a different result from one that fixes both. Both numbers
are already logged.

\notesec{Free, and worth doing on day one} Apply the post-hoc rescale by $r/s$ to the
existing best checkpoint. It changes no correlation and moves \file{nmse} from $1.010$ to
$0.944$. It is an evening's work and it removes the ``loses to the mean'' objection from
the figure regardless of what the campaign finds.

\subsection{Phase B, the pair term, under the winning objective\secstatus{todo}}

The reference operator provably cannot express any dependence on the pair (deleted gene,
reporter gene), and masked unmasking does not supply one at $k = 0$. The ladder is the
mechanism axis.

\begin{table}[htbp]
\centering\small
\begin{tabular}{llp{54mm}}
\toprule
arm & pair rank & note \\
\midrule
\texttt{P\_ref}      & 0  & additive reference \\
\texttt{P\_basis64}  & 64 & the current best-per-compute: $0.2274$ in 2.0 days \\
\texttt{P\_film}     & 90 & diagonal multiplicative, identity at init \\
\texttt{P\_graph}    & -- & masked graph attention applied \textbf{after} the perturbation operator, so the pair term is a function of network distance. \textbf{Not built.} \\
\bottomrule
\end{tabular}
\caption[]{Phase B. Four arms, two seeds: eight slots, $\approx 13$ GPU-days. \texttt{P\_graph}
requires new code and is gated on the graph-prior probe below.}
\label{tab:phase-b}
\end{table}

\notesec{Gate \texttt{P\_graph} on the probe, and run the probe this week} The probe
(\cref{sec:common}) costs minutes of CPU and asks whether deleting gene $X$ actually
perturbs $X$'s neighbors on each of the nine graphs. If the answer is no, \texttt{P\_graph}
is building a mechanism on a false premise and the nine attention heads currently spent
enforcing that premise should be freed instead. \textbf{Nothing in the campaign is more
decision-relevant per CPU-minute}, and it has been specified and unbuilt since 2026.07.26.

\subsection{Morphology runs alongside, and it is a different problem\secstatus{todo}}

Morphology is not a smaller copy of expression. Its target is 278-dimensional rather than
6,127, no morphology run has ever passed 121 epochs, and every one of the 397 runs across
eight projects trained on 1,161 strains (\cref{tab:morph-ntrain}) out of a 4,718-strain
screen.

\begin{enumerate}
\item \textbf{Drop \file{require_modalities} for the morphology-only arm} and rebuild the
  dataset on all 4,718 Ohya deletions. This is a query and config change, and it is the only
  way the $0.0824$ becomes a statement about morphology rather than about a quarter of it.
  \textbf{Do this first}; it is the largest single change available to the figure.
\item \textbf{Measure the per-epoch cost on the first run.} It is not known: the only
  figure on record is a tiny-model smoke at $\approx 47$ s/epoch against expression's
  $\approx 15$ in the same smoke, which would make morphology the more expensive of the two
  despite the smaller target. Until it is measured, the slot arithmetic above does not
  apply to morphology.
\item \textbf{Run the scalar warm-up in parallel}, on \file{A113} (actin\_n\_ratio, ceiling
  $0.873$, robust CV $2.37$) or \file{C124_C}. It calibrates the epoch budget for the
  278-feature target for a fraction of the cost.
\end{enumerate}

\subsection{The joint question, which is the one Figure 3 is about\secstatus{todo}}

Does expression help morphology. The control already exists and is correct:
\file{delta_joint_expr_morph_000.yaml} pins all three conditions to the 1,440 genotypes
carrying both phenotypes, so \texttt{expr\_morph} minus \texttt{morph} isolates the
auxiliary head from a data-quantity difference. It has never been run past 87 epochs.

\notesec{Run it at two scales, and do not conflate them} On the 1,440 shared strains it is
the mechanism test it was designed to be. Paired with \textbf{fitness} instead of
expression it keeps 4,220 of the 4,718 strains, three times the instance count, which is
where the statistical power is. The two answer different questions and both are worth a
slot; the fitness pairing is the one likely to produce a detectable effect.

\notesec{The prior from the metabolite strands is not encouraging, and should be said out
loud} The one auxiliary-head experiment run to completion anywhere in this project is the
metabolome head on betaxanthin, and it reads $-0.0265 \pm 0.0159$
(\cref{tab:bx-aa-paired}). A second phenotype sharing an encoder has not yet been shown to
help any first phenotype here. The expression-morphology pairing has a better prior than
that one, because the two phenotypes share strains rather than merely genes, but the
campaign should be designed to \textbf{measure} the effect rather than to confirm it, which
means the control arm matters more than the joint arm.

\subsection{What the campaign cannot fix\secstatus{todo}}

Expression sits at $0.238$ against a ceiling of $0.775$, and the perturbation axis has 1,484
points with each gene deleted exactly once (\cref{tab:label-accounting}). No objective and
no pair term changes that. If Phase A and Phase B both come back inside the noise floor,
the honest reading is that the binding constraint is the data rather than the model, and
the figure's expression panel should be built around the ceiling-relative statement and the
imputation capability rather than around a genotype-to-expression number.
